\section{局部紧致性}\label{section:ch03-locally-compact}

本节我们研究局部紧致性这个概念，并且证明一个基本定理：任何一个局部紧致的 Hausdorff 空间可以嵌入到某一个叫做单点紧致化的紧致的 Hausdorff 空间。

\begin{definition}\label{definition:ch03-locally-compact}
空间 $X$ 是在 $x$ 处\textbf{局部紧致的}（locally compact at $x$），若存在 $X$ 的一个紧致子空间 $C$ 包含着 $x$ 的一个邻域。如果 $X$ 在它的每一点处都是局部紧致的，则称 $X$ 是\textbf{局部紧致的}（locally compact）。
\end{definition}

一个紧致空间自然是局部紧致的。

\begin{example}\label{example:ch03-real-locally-compact}
实直线 $\mathbb{R}$ 是局部紧致的。因为点 $x$ 属于某一个区间 $(a,b)$，而 $(a,b)$ 又包含在紧致空间 $[a,b]$ 中。读者可以验证，$\mathbb{R}$ 的有理数子空间 $\mathbb{Q}$ 不是局部紧致的。
\end{example}

\begin{example}\label{example:ch03-rn-locally-compact}
空间 $\mathbb{R}^{n}$ 是局部紧致的。点 $x$ 属于某一个基元素 $(a_{1},b_{1})\times\cdots\times(a_{n},b_{n})$，而这个基元素又包含在紧致空间 $[a_{1},b_{1}]\times\cdots\times[a_{n},b_{n}]$ 之中。空间 $\mathbb{R}^{\omega}$ 不是局部紧致的。因为它的任何基元素都不被紧致子空间所包含。事实上，若
\begin{align*}
B=(a_{1},b_{1})\times\cdots\times(a_{n},b_{n})\times\mathbb{R}\times\cdots\times\mathbb{R}\times\cdots
\end{align*}
被某一个紧致子空间所包含，那么它的闭包
\begin{align*}
\overline{B}=[a_{1},b_{1}]\times\cdots\times[a_{n},b_{n}]\times\mathbb{R}\times\cdots
\end{align*}
就是紧致的，但它并不紧致。
\end{example}

\begin{example}\label{example:ch03-lots-locally-compact}
任何一个具有上确界性质的全序集 $X$ 是局部紧致的：任意给定 $X$ 的一个基元素，它一定包含在 $X$ 的一个闭区间中，而闭区间是紧致的。
\end{example}

数学中所讨论的两类属性最好的空间当属可度量化空间和紧致的 Hausdorff 空间。它们有许多有用的性质，可用于证明定理、建立构架等等。如果所给的空间不是上述两类空间之一，则退而求其次，我们希望它是这两类空间之一的子空间。当然，可度量化空间的子空间还是可度量化的，因而这种方式不能得到任何新空间，然而，紧致的 Hausdorff 空间的子空间就未必紧致了。于是就有这样一个问题：在什么条件下，一个空间同胚于紧致的 Hausdorff 空间的一个子空间？这里先给出一个回答。在第 5 章研究了一般的紧致化问题之后，我们再来讨论这个问题。

\begin{theorem}\label{theorem:ch03-one-point-compactification}
设 $X$ 是一个空间。则 $X$ 是一个局部紧致的 Hausdorff 空间当且仅当存在一个空间 $Y$ 使得以下条件成立：
\begin{enumerate}
\item $X$ 是 $Y$ 的子空间.
\item 集合 $Y\setminus X$ 是单点集.
\item $Y$ 是紧致的 Hausdorff 空间.
\end{enumerate}
若 $Y$ 和 $Y'$ 是满足上述条件的两个空间，则存在从 $Y$ 到 $Y'$ 的一个同胚使得它在 $X$ 上的限制是恒等映射。
\end{theorem}

\begin{proof}
第一步。我们首先验证唯一性。设 $Y$ 和 $Y'$ 是满足上述条件的两个空间。定义映射 $h\colon Y\to Y'$ 使得 $h$ 将 $Y\setminus X$ 的那个单点 $p$ 映为 $Y'\setminus X$ 的那个单点 $q$；$h$ 在 $X$ 上是恒等映射。我们证明：若 $U$ 是 $Y$ 的开集，则 $h(U)$ 是 $Y'$ 的一个开集。那么，对称性便蕴涵着 $h$ 是一个同胚。

首先，考虑 $U$ 不包含 $p$ 的情形。这时 $h(U)=U$。由于 $U$ 是 $Y$ 的一个开集，并且它包含于 $X$，从而 $U$ 是 $X$ 的一个开集。由于 $X$ 是 $Y'$ 的一个开集，因此 $U$ 也是 $Y'$ 的一个开集。

其次，考虑 $U$ 包含着 $p$ 的情形。由于 $C=Y\setminus U$ 是 $Y$ 的一个闭集，因此 $C$ 作为 $Y$ 的子空间是紧致的。因为 $C$ 包含于 $X$，从而它是 $X$ 的一个紧致子空间。由于 $X$ 是 $Y'$ 的子空间，因此空间 $C$ 也是 $Y'$ 的紧致子空间。由于 $Y'$ 是一个 Hausdorff 空间，因此 $C$ 是 $Y'$ 的一个闭集，从而 $h(U)=Y'\setminus C$ 是 $Y'$ 的一个开集。

第二步。设 $X$ 是一个局部紧致的 Hausdorff 空间，我们来构造空间 $Y$。第一步的讨论已为我们提供了构造 $Y$ 的思路。为了方便起见，用符号 $\infty$ 表示不属于 $X$ 的某一个元素，并且把它添加到 $X$ 上构成集合 $Y=X\cup\{\infty\}$，构成 $Y$ 的拓扑的 $Y$ 的开集族取为所有下列类型的集合：类型一：$X$ 的开子集 $U$。类型二：$Y\setminus C$，其中 $C$ 是 $X$ 的一个紧致子空间。

我们需要验证这个族的确是 $Y$ 的一个拓扑。空集是属于类型一的，空间 $Y$ 是属于类型二的。验证两个开集的交是开集可考虑下面三种情况：
\begin{align*}
U_{1}\cap U_{2}\quad\qquad\qquad\qquad\text{是类型一的.}
\end{align*}
\begin{align*}
(Y\setminus C_{1})\cap(Y\setminus C_{2})=Y\setminus(C_{1}\cup C_{2})\quad\text{是类型二的.}
\end{align*}
\begin{align*}
U_{1}\cap(Y\setminus C_{1})=U_{1}\cap(X\setminus C_{1})\quad\text{是类型一的.}
\end{align*}
最后一条是因为 $C_{1}$ 在 $X$ 中是闭的。同样可以证明开集族的任意并是开的：
\begin{align*}
\bigcup U_{\alpha}=U\quad\qquad\qquad\qquad\text{是类型一的.}
\end{align*}
\begin{align*}
\bigcup(Y\setminus C_{\beta})=Y\setminus(\cap C_{\beta})=Y\setminus C\quad\text{是类型二的.}
\end{align*}
\begin{align*}
(\bigcup U_{\alpha})\cup(\bigcup(Y\setminus C_{\beta}))=U\cup(Y\setminus C)=Y\setminus(C-U)
\end{align*}
是类型二的，这是因为 $C-U$ 是 $C$ 的闭子集，因而是紧致的。

以下证明 $X$ 是 $Y$ 的一个子空间。任意给定 $Y$ 的开集，先证明它与 $X$ 的交是 $X$ 的开集。这是因为：若 $U$ 为类型一的集合，则 $U\cap X=U$。若 $Y\setminus C$ 为类型二的集合，则 $(Y\setminus C)\cap X=X\setminus C$，两者都是 $X$ 中的开集。反之，$X$ 中任何开集都是类型一的集合，因而，据定义得知它也是 $Y$ 的开集。

为了证明 $Y$ 是紧致的，设 $\mathcal{A}$ 是 $Y$ 的一个开覆盖。族 $\mathcal{A}$ 一定包含一个类型二的开集，比如说 $Y\setminus C$，这是因为在类型一中没有包含着点 $\infty$ 的开集。把 $\mathcal{A}$ 中所有那些不同于 $Y\setminus C$ 的集合取出来与 $X$ 作交，它们是由 $X$ 中开集所组成的 $C$ 覆盖。因为 $C$ 是紧致的，所以它们中的有限多个可以覆盖 $C$，则 $\mathcal{A}$ 中与之相应的成员所构成的有限族连同 $Y\setminus C$ 便覆盖着 $Y$。

为了证明 $Y$ 是 Hausdorff 的，设 $x$ 和 $y$ 是 $Y$ 的两点。若它们都属于 $X$，则存在 $X$ 中无交的开集 $U$ 和 $V$ 分别包含它们。否则，不妨设 $x\in X$ 和 $y=\infty$，那么可以在 $X$ 中选择包含 $x$ 的邻域 $U$ 的一个紧致集 $C$。于是 $U$ 和 $Y\setminus C$ 便是 $Y$ 中分别包含 $x$ 和 $\infty$ 的两个无交的邻域。

第三步。最后，我们证明充分性。假设满足条件 (1)--(3) 的空间 $Y$ 存在。由于 $X$ 是一个 Hausdorff 空间 $Y$ 的一个子空间，所以 $X$ 是一个 Hausdorff 空间。任意给定 $x\in X$，我们证明 $X$ 在点 $x$ 处是局部紧致的。在 $Y$ 中分别选取包含 $x$ 和单点 $Y\setminus X$ 的无交开集 $U$ 和 $V$。那么 $C=Y\setminus V$ 是 $Y$ 的一个闭集，从而是 $Y$ 的紧致子空间。由于 $C$ 包含于 $X$，它也是 $X$ 的紧致子空间，并且它包含着 $x$ 的一个邻域 $U$。
\end{proof}

在上述定理中，如果 $X$ 自身就是一个紧致空间，则构造空间 $Y$ 的意义就不大，因为此时 $Y$ 是 $X$ 加上一个孤立点。但当 $X$ 不紧致时，那么 $Y\setminus X$ 中的那个单点是 $X$ 的一个极限点，从而，$\overline{X}=Y$。

\begin{definition}\label{definition:ch03-compactification}
若 $Y$ 是一个紧致的 Hausdorff 空间，$X$ 是 $Y$ 的真子空间并且其闭包等于 $Y$，则 $Y$ 称为 $X$ 的一个\textbf{紧致化}（compactification）。若 $Y\setminus X$ 为单点集，则 $Y$ 称为 $X$ 的\textbf{单点紧致化}（one-point compactification）。
\end{definition}

前述结论表明：$X$ 具有单点紧致化 $Y$ 当且仅当 $X$ 为非紧致的局部紧致的 Hausdorff 空间。之所以不将 $Y$ 称为``一个''单点紧致化是由于在同胚的意义下 $Y$ 是唯一的。

\begin{example}\label{example:ch03-one-point-compact-examples}
读者可以验证，实直线 $\mathbb{R}$ 的单点紧致化同胚于圆周。同样，$\mathbb{R}^{2}$ 的单点紧致化同胚于球面 $S^{2}$。如果将 $\mathbb{R}^{2}$ 看成复数空间 $\mathbb{C}$，则 $\mathbb{C}\cup\{\infty\}$ 称为 Riemann 球面或扩充复平面。
\end{example}

我们所给出的这个局部紧致性的定义，有些地方还不能令人十分满意。通常说一个空间 $X$``局部地''满足某一个性质，是指对于每一个 $x\in X$ 都存在一个``任意小''的邻域具有所说的性质。而这个局部紧致的定义并没有涉及到``任意小''的邻域，于是便产生了它到底该不该称为局部紧致性的问题。

下面是局部紧致性的另一种说法，它比较自然地体现了对``局部''的要求，而且当 $X$ 是一个 Hausdorff 空间时，它等价于我们的定义。

\begin{theorem}\label{theorem:ch03-locally-compact-neighborhood}
设 $X$ 是一个 Hausdorff 空间。则 $X$ 在 $x$ 处局部紧致当且仅当对于 $x$ 的任何一个邻域 $U$，存在 $x$ 的一个邻域 $V$，使得 $\overline{V}$ 紧致并且 $\overline{V}\subset U$。
\end{theorem}

\begin{proof}
显然，这个新的表述蕴涵着局部紧致性。$C=\overline{V}$ 便是所要求的包含 $x$ 的一个邻域的紧致集。反之，假设 $X$ 是局部紧致的，$x$ 是 $X$ 的一个点，$U$ 是 $x$ 的任何一个邻域。记 $Y$ 为 $X$ 的单点紧致化，集合 $C$ 为 $Y\setminus U$。则 $C$ 是 $Y$ 的一个闭子集，从而 $C$ 是 $Y$ 的一个紧致子空间。应用\cref{lemma:ch03-separate-point-and-compact}，选取分别包含 $x$ 和 $C$ 的无交开集 $V$ 和 $W$。于是 $V$ 的闭包 $\overline{V}$ 是 $Y$ 的一个紧致子空间，$\overline{V}$ 与 $C$ 无交，从而 $\overline{V}\subset U$。
\end{proof}

\begin{corollary}\label{corollary:ch03-locally-compact-subspace}
设 $X$ 是局部紧致的 Hausdorff 空间，$A$ 是 $X$ 的一个子空间。若 $A$ 是 $X$ 的一个闭集或者一个开集，则 $A$ 是局部紧致的。
\end{corollary}

\begin{proof}
假设 $A$ 是 $X$ 的一个闭集。对于 $x\in A$，令 $C$ 为 $X$ 中的一个紧致子空间，并且包含着 $x$ 在 $X$ 中的一个邻域 $U$。那么 $C\cap A$ 是 $C$ 的一个闭集，因而是紧致的，并且它包含着 $x$ 在 $A$ 中的邻域 $U\cap A$。（这里还没有用到 Hausdorff 条件。）

假设 $A$ 是 $X$ 中的一个开集，对于 $x\in A$，我们应用上述定理，选取 $x$ 在 $X$ 中的一个邻域 $V$，使得 $\overline{V}$ 是紧致的并且 $\overline{V}\subset A$。那么 $C=\overline{V}$ 是 $A$ 的一个紧致子空间，它包含着 $x$ 在 $A$ 中的邻域 $V$。
\end{proof}

\begin{corollary}\label{corollary:ch03-locally-compact-open-subspace}
空间 $X$ 同胚于一个紧致的 Hausdorff 空间的开子集当且仅当 $X$ 是一个局部紧致的 Hausdorff 空间。
\end{corollary}

\begin{proof}
据\cref{theorem:ch03-one-point-compactification} 及\cref{corollary:ch03-locally-compact-subspace} 便可得到结论。
\end{proof}

\subsection*{习题}

\begin{enumerate}
\item 证明：有理数集 $\mathbb{Q}$ 不是局部紧致的。
\item 设 $\{X_{\alpha}\}$ 为非空空间的一个加标族。
\begin{enumerate}
\item 证明：若 $\prod X_{\alpha}$ 是局部紧致的，则每一个 $X_{\alpha}$ 是局部紧致的并且除去有限多个 $\alpha$ 之外，$X_{\alpha}$ 是紧致的。
\item 应用 Tychonoff 定理的结论，证明 (a) 的逆命题。
\end{enumerate}
\item 设 $X$ 是一个局部紧致空间。$f\colon X\to Y$ 连续是否意味着空间 $f(X)$ 是局部紧致的？$f$ 是连续开映射时结论如何？验证你的结论。
\item 证明：$[0,1]^{\omega}$ 关于一致拓扑不是局部紧致的。
\item 如果 $f\colon X_{1}\to X_{2}$ 是局部紧致的 Hausdorff 空间之间的一个同胚，证明 $f$ 可扩充为它们的单点紧致化之间的同胚。
\item 证明：$\mathbb{R}$ 的单点紧致化同胚于圆周 $S^{1}$。
\item 证明：$S_{\Omega}$ 的单点紧致化同胚于 $\overline{S}_{\Omega}$。
\item 证明：$\mathbb{Z}_{+}$ 的单点紧致化同胚于 $\mathbb{R}$ 的子空间 $\{0\}\cup\{1/n\mid n\in\mathbb{Z}_{+}\}$。
\item 证明：若 $G$ 是一个局部紧致的拓扑群，$H$ 是它的一个子群，则 $G/H$ 是局部紧致的。
\item 设 $X$ 是一个 Hausdorff 空间，并且在点 $x$ 处是局部紧致的。证明：对于 $x$ 的任意邻域 $U$，存在 $x$ 的一个邻域 $V$，使得 $\overline{V}$ 是紧致的，并且 $\overline{V}\subset U$。
\item 证明以下结论。
\begin{enumerate}
\item \textbf{引理}\quad 设 $p\colon X\to Y$ 是一个商映射并且 $Z$ 是一个局部紧致的 Hausdorff 空间，则映射
\begin{align*}
\pi=p\times i_{Z}\colon X\times Z\longrightarrow Y\times Z
\end{align*}
是一个商映射。[提示：若 $\pi^{-1}(A)$ 是包含 $x\times y$ 的一个开集，选取开集 $U_{1}$ 和 $V$（其中 $\overline{V}$ 是紧致的），使得 $x\times y\in U_{1}\times V$ 和 $U_{1}\times\overline{V}\subset\pi^{-1}(A)$。对于 $U_{i}\times\overline{V}\subset\pi^{-1}(A)$，应用管状引理，选取包含 $p^{-1}(p(U_{i}))$ 的一个开集 $U_{i+1}$，使得 $U_{i+1}\times\overline{V}\subset\pi^{-1}(A)$。令 $U=\bigcup U_{i}$。证明 $U\times V$ 是 $x\times y$ 的一个饱和邻域，并且包含在 $\pi^{-1}(A)$ 中。]

在第 46 节的习题中将给出一个与上述完全不同的证明方法的概要。
\item \textbf{定理}\quad 设 $p\colon A\to B$ 和 $q\colon C\to D$ 都是商映射。若 $B$ 和 $C$ 是局部紧致的 Hausdorff 空间，则 $p\times q\colon A\times C\to B\times D$ 是一个商映射。
\end{enumerate}
\end{enumerate}